% Preamble for: An Explicit Exponential-Period Representation of the V_quad Connection Coefficient
\documentclass[11pt,reqno]{amsart}
\usepackage[T1]{fontenc}
\usepackage{amsmath,amssymb,amsthm,mathtools}
\usepackage{microtype}
\usepackage[letterpaper,margin=1.5in]{geometry}
\usepackage{enumitem}
\usepackage{booktabs}
\usepackage{tikz}
\usetikzlibrary{arrows.meta,decorations.markings,decorations.pathmorphing}
\usepackage[colorlinks=true,linkcolor=blue,citecolor=blue,urlcolor=blue]{hyperref}

% --- byte-reproducible build guards (memory: reproducible PDFs) ---
\pdfinfoomitdate=1
\pdftrailerid{}
\pdfsuppressptexinfo=-1

% --- theorem environments ---
\theoremstyle{plain}
\newtheorem{theorem}{Theorem}[section]
\newtheorem{proposition}[theorem]{Proposition}
\newtheorem{lemma}[theorem]{Lemma}
\newtheorem{corollary}[theorem]{Corollary}
\theoremstyle{definition}
\newtheorem{definition}[theorem]{Definition}
\newtheorem{remark}[theorem]{Remark}
\newtheorem{example}[theorem]{Example}

% --- macros ---
\newcommand{\Vquad}{V_{\mathrm{quad}}}
\newcommand{\Bhat}{\widehat{B}}
\newcommand{\Lphi}{L_{\varphi}}
\newcommand{\LV}{L_{V}}
\newcommand{\Qsqrt}{\mathbb{Q}(\sqrt3)}
\newcommand{\Qbar}{\overline{\mathbb{Q}}}
\newcommand{\xizero}{\xi_0}
\newcommand{\SL}{\mathrm{SL}}
\newcommand{\GL}{\mathrm{GL}}
\newcommand{\Gm}{\mathbb{G}_m}
\providecommand{\C}{}\renewcommand{\C}{C}
\DeclareMathOperator*{\K}{K}

\title[An exponential-period representation of the $V_{\mathrm{quad}}$ constant]%
{An explicit exponential-period representation of the $V_{\mathrm{quad}}$ connection coefficient}
\author{Papanokechi}
\thanks{ORCID \texttt{0009-0000-6192-8273}.}
\date{15 June 2026}

\begin{document}
\begin{abstract}
We prove that the connection coefficient $C$ of the $\Vquad$ polynomial continued fraction---the
rank-one standalone closure on the Sakai surface $D_5^{(1)}$ (Painlev\'e~V)---admits an explicit
exponential-period representation
$C=(\lvert\Gamma(\beta)\rvert/2\pi)\int_\gamma e^{\xi}\Bhat(\xi)\,d\xi=\lvert\Gamma(\beta)\rvert K$,
where $\Bhat$ is the Borel transform of the $\Vquad$ asymptotic series, holonomic of order $4$
(the Borel--Laplace dual of the order-$2$ operator annihilating the series)
with coefficients in the real quadratic field $\Qsqrt$; $\gamma$ is an explicit Hankel
rapid-decay cycle on $(-\infty,-2/\sqrt3]$; and $\beta=-1/(3\sqrt3)$. The identity is verified by
three structurally independent methods---differential-equation/operator duality, Borel--Laplace
contour deformation, and Stokes-data---agreeing to $46$ digits. The order-$2$ operator annihilating
the asymptotic series has differential Galois group $\SL_2(\mathbb{C})$ by Kovacic's algorithm.
As a by-product, holonomicity of the Borel transform forces a finite resurgent structure. Finally,
conditional on the Fres\'an--Jossen period conjecture for exponential motives and on a stated
motivic-comparison hypothesis, $C$ is transcendental over $\Qbar$.
\end{abstract}
\maketitle
